\begin{ouverture}[calcul]{6}{Additionner et soustraire des fractions}{Partie 1 · Nombres et calculs}
{Dénominateurs quelconques · Quantités · Partages}
\objectif{Additionner et soustraire avec un dénominateur commun.}
\objectif{Résoudre des problèmes de parts et de restes.}
\objectif{Contrôler une somme par rapport à l’unité.}
\prerequis{Fractions égales et multiples.}
\end{ouverture}
\pagestyle{otMain}\essentiel\fontsize{11}{15}\selectfont
\phantomsection\label{cours-6-1}
\essSection{1}{Exprimer les parts de la même façon}
\begin{propriete}Pour $b>0$, $\frac ab+\frac cb=\frac{a+c}b$ et $\frac ab-\frac cb=\frac{a-c}b$. Le dénominateur indique la taille des parts ; il reste inchangé.\end{propriete}
\label{fig-cours-6-1-cours-somme-fractions}\figureGeo{cours-somme-fractions}

\begin{savoirfaire}[Changer de partage]Pour $\frac34+\frac56$, utiliser le dénominateur 12 : $\frac9{12}+\frac{10}{12}=\frac{19}{12}=1+\frac7{12}$. Un multiple commun convient ; prendre un petit multiple réduit les calculs.\end{savoirfaire}
\begin{attention}$\frac12+\frac13$ ne vaut pas $\frac25$. La somme de $1/2$ et d’un nombre positif est supérieure à $1/2$, alors que $2/5<1/2$ : ce contrôle révèle l’erreur.\end{attention}
\clearpage
\phantomsection\label{cours-6-2}
\essSection{2}{Interpréter les fractions}
\begin{savoirfaire}[Fraction d’une quantité]Prendre $a/b$ d’une quantité signifie partager cette quantité en $b$ parts égales puis en prendre $a$. Les $3/5$ de 40 valent $(40:5)\times3=24$. On peut raisonner en parts sans utiliser de règle de produit de fractions.\end{savoirfaire}
\begin{exemple}Un trajet est parcouru pour $2/5$ le matin et $1/3$ l’après-midi, ces deux fractions portant sur le trajet entier. La part faite est $6/15+5/15=11/15$ ; il reste $1-11/15=4/15$.\end{exemple}
\begin{savoirfaire}[Identifier le tout]Lorsqu’on parle d’une fraction du reste, calculer le reste avant d’appliquer la fraction. Pour 60 L, après usage de $1/3$ du total il reste 40 L ; la moitié de ce reste vaut 20 L.\end{savoirfaire}
\begin{attention}Les questions de ce chapitre portent sur des fractions positives et des additions ou soustractions. Un partage peut donner une quantité non entière ; vérifier si le contexte l’autorise.\end{attention}

\clearpage\rubrique{Découverte}
\begin{blocLivre}{ex-6-1}\begin{exercice}[1]{Addition directe}
Calculer $2/7+3/7$.
\end{exercice}\end{blocLivre}
\begin{blocLivre}{ex-6-2}\begin{exercice}[1]{Soustraction directe}
Calculer $9/11-4/11$.
\end{exercice}\end{blocLivre}
\begin{blocLivre}{ex-6-3}\begin{exercice}[1]{Complément à un}
Calculer $1-3/8$.
\end{exercice}\end{blocLivre}
\begin{blocLivre}{ex-6-4}\begin{exercice}[1]{Dénominateurs multiples}
Calculer $1/3+1/6$.
\end{exercice}\end{blocLivre}
\begin{blocLivre}{ex-6-5}\begin{exercice}[1]{Une part de 28}
Calculer les $3/4$ de 28.
\end{exercice}\end{blocLivre}

\rubrique{Entraînement}
\begin{blocLivre}{ex-6-6}\begin{exercice}[2]{Dénominateurs quelconques}
Calculer $5/8+7/12$.
\end{exercice}\end{blocLivre}
\begin{blocLivre}{ex-6-7}\begin{exercice}[2]{Une différence}
Calculer $7/10-5/12$.
\end{exercice}\end{blocLivre}
\begin{blocLivre}{ex-6-8}\begin{exercice}[2]{Trois termes}
Calculer $2/3+3/4-5/6$.
\end{exercice}\end{blocLivre}
\begin{blocLivre}{ex-6-9}\begin{exercice}[2]{Entier et fraction}
Calculer $3+5/6-7/4$.
\end{exercice}\end{blocLivre}
\begin{blocLivre}{ex-6-10}\begin{exercice}[2]{Retrouver une part}
Compléter $5/9+\,?=7/6$.
\end{exercice}\end{blocLivre}

\rubrique{Problèmes}
\begin{blocLivre}{ex-6-11}\begin{exercice}[3]{Deux étapes de trajet}
On parcourt $3/8$ d’un trajet le matin et $5/12$ du trajet entier l’après-midi. Quelle fraction reste-t-il ? Pour 144 km, quelle distance est-ce ?
\end{exercice}\end{blocLivre}
\begin{blocLivre}{ex-6-12}\begin{exercice}[3]{Un budget}
Sur un budget de 240 euros, $1/3$ sert aux transports et $3/8$ aux repas. Le reste suffit-il pour un achat de 75 euros ?
\end{exercice}\end{blocLivre}
\begin{blocLivre}{ex-6-13}\begin{exercice}[3]{Fraction du reste}
Une citerne contient 90 L. On utilise $2/5$ du contenu, puis $1/3$ du reste. Combien reste-t-il après les deux usages ?
\end{exercice}\end{blocLivre}
\begin{blocLivre}{ex-6-14}\begin{exercice}[3]{Retrouver le total}
Un ruban a été coupé : $1/4$ de sa longueur pour A, $1/3$ pour B, tous deux du ruban initial. Les 35 cm restants sont pour C. Trouver la longueur initiale.
\end{exercice}\end{blocLivre}
\begin{blocLivre}{ex-6-15}\begin{exercice}[3]{Un partage impossible}
Une élève promet $5/8$ d’un gâteau à un groupe et $7/12$ du même gâteau à un autre, sans chevauchement. Est-ce possible ? De quelle fraction la promesse dépasse-t-elle le gâteau ?
\end{exercice}\end{blocLivre}

\clearpage
